\scanpage{3}\mspage{2}

prenant pour $A$ la ``champ'' sur $X$ formé des $F$-torseurs à droite, la (Gr)- catégorie des au\deleted{ot}to-équivalences de $A$ avec lui-même s'interprète comme la catégorie des ``b\deleted{t}itorseurs'' sous $F$, i.e. des faisceaux \inserted{\handwritten{$P$}} sur lesquels $F$ opère à la fois à droite et à gauche, ces opérations commutant et chacune d'elles faisant de $P$ un torseur (à droite ou à gauche) sous $F$, - la composition $\otimes$ étant la composition de Baer évident\deleted{\unsure{e|s}}e. Lorsque $F$ est encore de la forme $\underline{O}_{X}^{*}$ ($\underline{O}_{X}$ étant un' faisceau d'anneaux, qu'on ne suppose plus nécessairement commutatif) ces bitorseurs s'interprètent aussi en termes de bi-Modules \handwritten{``}inversibles\handwritten{''} sous $\underline{O}_{X}$.

\uline{Structure}. \deleted{\rlap{xxx}\unsure{Les|les}} Soit $C$ une (Gr)-catégorie, on lui associe

a) le groupe $\handwritten{\pi}_{0}(C)$ des classes d'isomorphisme d'objets de $C$,

b) le groupe $\handwritten{\pi}_{1}(C)$ des automorphismes de $1_C$ (objet unité de $C$)

c) une action de $\handwritten{\pi}_{0}(C)$ sur $\handwritten{\pi}_{1}(C)$, en associant à tout objet $X$ de $C$ l'automorphisme \deleted{\rlap{xxxx}de 1} $\handwritten{\rho}(X)$ de $1_C$ déduit des deux isomorphismes
\diagram{\begin{tikzcd}[ampersand replacement=\&]
\mathrm{Aut}(1_C) \handwritten{\arrow[r,shift left=.55ex]\arrow[r,shift right=.55ex]} \& \mathrm{Aut}(X)
\end{tikzcd}}
donnés par $u\handwritten{\mapsto}u\otimes\mathrm{id}_{X}$ et $u\handwritten{\mapsto}\mathrm{id}_{X}\otimes u$. On prouve que $\handwritten{\pi}_{1}(C)$ est un groupe \uline{commutatif} et que l'on obtient bien par c) une structure de $\handwritten{\pi}_{0}(C)$-module sur celui-ci. Ceci posé, si on choisit pour tout $a\handwritten{\in}\handwritten{\pi}_{0}(C)$ un représentant $L_{a}\handwritten{\in}\mathrm{Ob}\ C$ de $a$, et pour deux \deleted{\rlap{xxxxxx}objets} $a,b$ un isomorphisme
\[
\handwritten{\phi}_{a,b}: L_a\otimes L_b\handwritten{\simeq}L_{ab}\ ,
\]
alors pour trois éléments $a,b,c\handwritten{\in}\handwritten{\pi}_{0}(C)$, l'isomorphisme d'associativité
\[
(L_a\otimes L_b)\otimes L_c\handwritten{\simeq}L_a\otimes(L_b\otimes L_c)\ ,
\]
compte tenu des isomorphismes $\handwritten{\phi}_{a,b}$, $\handwritten{\phi}_{ab,c}$, $\handwritten{\phi}_{b,c}$ et $\handwritten{\phi}_{a,bc}$, peut s'interpréter comme un isomorphisme
\[
L_{abc}\handwritten{\simeq}L_{abc}\ ,
\]
ou encore comme la tensorisation \deleted{\unsure{à|a}}à gauche avec un élément bien déterminé
\[
f(a,b,c)\handwritten{\in}\handwritten{\pi}_{1}(C)\deleted{\rlap{$\mathrm{xxxxxxxxx}_{\mathrm{xxxxxx}}$}1_{C}\quad 1_{C}}\ .
\]
% No connector is recoverable between the two surviving cancelled units.
Les $\handwritten{\phi}_{a,b}$ étant choisis, on voit donc que la donnée d'un isomorphisme d'associativité fonctoriel $(L\otimes L')\otimes L''\handwritten{\simeq}L\otimes(L'\otimes L'')$ équivaut à la donnée de l'application
\[
f:\handwritten{\pi}_{0}\mathbin{\mathrm{x}}\handwritten{\pi}_{0}\mathbin{\mathrm{x}}\handwritten{\pi}_{0}\handwritten{\longrightarrow}\handwritten{\pi}_{1}\ .
\]
On vérifie alors que l'axiome \deleted{\handwritten{\unsure{(du pen}}} standard d'autocompatibilité d'une donnée d'associativité \inserted{\handwritten{(axiome du pentagone)}} s'exprime précisément par la condition que $f$ soit un \handwritten{3}-\uline{cocycle} du groupe $\handwritten{\pi}_{\deleted{1}\handwritten{0}}$ à valeurs dans le groupe $\handwritten{\pi}_{1}$. D'autre part, l'indétermination dans le choix du système d'isomorphismes $\handwritten{\phi}_{a,b}$ est précisément donnée par une 2-cochaine arbitraire, et on voit que si on change $\handwritten{\phi}$ par une 2-coc\deleted{a}h\deleted{h}aine $g$, $f$ est changé en $f+dg$ - donc l'ensemble des $f$ \deleted{p}co\deleted{s}rrespondants à des choix différents de $\handwritten{\phi}$ est exactement une classe de \handwritten{3}-cohomologie
% archive sheet number: 2
